\section{Harness AlignmentGraph-ISD}
\label{sec:method}

\subsection{Tổng quan}
\label{sec:overview}

Bài toán ở \cref{sec:problem-statement} đòi hỏi alignment phải được sinh ra và đảm
bảo độc lập với model. Quan sát dẫn dắt thiết kế của chúng tôi là: mọi thứ điều này
cần---trạng thái bền vững, kiểm soát chất lượng, và điều phối chuyên biệt---là trách
nhiệm của môi trường thực thi, không phải năng lực của model. Do đó AlignmentGraph-ISD
tập trung các trách nhiệm này vào một harness, để model chỉ sinh nội dung bên trong
nó. Một agent hoàn chỉnh là sự kết hợp của cả hai (\emph{agent = harness + model}).
Cụ thể, harness (i)~phân rã việc sinh cho các agent chuyên biệt xây dựng alignment
graph theo một thứ tự thiết kế giảng dạy cố định với phạm vi ghi giới hạn
(\cref{sec:decomposition}), để alignment hình thành by construction, và (ii)~trao
cho mỗi agent pedagogical verification cùng repair như các công cụ nó chạy ngay
trong lượt của mình (\cref{sec:qc}), để vi phạm được bắt và sửa trước khi việc sinh
tiếp tục. Một bước aggregation cuối (\cref{sec:aggregation}) phát ra chương trình và
ghi lại một residual audit.

\subsection{Multi-Agent Decomposition trên Graph}
\label{sec:decomposition}

Mỗi concern ISD do một specialist agent xử lý---needs context, objective design,
assessment design, activity và material, implementation plan, và evaluation---chạy
theo thứ tự ADDIE cố định. Mỗi agent $A_i$ là một toán tử cập-nhật-graph có phạm vi
ghi bị giới hạn bởi chuyên môn của nó:
\[
A_i : (G_t, c) \mapsto G_{t+1}, \qquad \Delta(G_t, G_{t+1}) \subseteq V^{(i)} \cup E^{(i)},
\]
trong đó $\Delta(G_t, G_{t+1})$ là các node và edge được thêm giữa hai graph, còn
$V^{(i)}, E^{(i)}$ là các kiểu node và edge mà agent $i$ được ghi. Ví dụ, assessment
agent chỉ ghi node \texttt{assessment} và edge \textsl{measures}: khi nó phát ra một
assessment item nó phải nêu tên objective mà item chấm, và việc nêu tên đó \emph{chính
là} edge \textsl{measures}. Vì objective agent chạy trước và đặt các objective vào
graph chung, mọi agent sau tham chiếu tới objective đã có, và các edge alignment của
\cref{eq:aligned} được tạo ra như một phần bắt buộc của việc sinh. Nhờ vậy alignment
đúng \emph{by construction}, thay vì được kỳ vọng trong văn bản. Context mỗi agent
nhận là $c = \rho(b) \cup \gamma(G_t) \cup \mu(M)$, với $\rho(b)$ là context từ
brief, $\gamma(G_t)$ là graph-context đọc từ các node trước, và $\mu(M)$ là bộ nhớ
domain ổn định (ADDIE, Bloom, rubric template).

\subsection{Verification và Repair như Công cụ Per-Agent}
\label{sec:qc}

Chỉ riêng construction đảm bảo các edge \emph{tồn tại}; nó không đảm bảo mỗi edge
\emph{hợp lệ} hay mọi objective được phủ. Đó là việc của verification và repair.
Quan trọng: đây không phải một giai đoạn tổng thể riêng chạy sau khi sinh xong; chúng
là các công cụ mỗi specialist chạy \emph{ngay trong lượt của mình}, ngay sau khi viết
artifact. Điều này giữ cho trajectory của mỗi agent là một vòng sạch
generate\,$\rightarrow$\,self-check\,$\rightarrow$\,emit và cho phép các kiểm tra quan
hệ hoạt động, vì đến lúc một agent self-check, các artifact trước mà nó phụ thuộc đã
có trong graph chung.

\paragraph{Verifier.} Gọi $\mathcal{V}$ là pedagogical verifier và
$\Omega = \mathcal{V}(G)$ là tập vi phạm nó phát hiện. Nó chạy bảy họ ràng buộc tất
định trên graph: objective coverage (sự tồn tại
\textsl{measures}/\textsl{supports}/\textsl{evaluates} theo từng objective), tính
hợp lệ của Bloom-verb và nhất quán verb--level~\cite{anderson2001taxonomy}, khớp
\emph{cấp độ nhận thức} assessment--objective (một objective bậc cao chỉ đo bằng item
mức nhận diện sẽ bị đánh dấu), nhất quán rubric (tham chiếu objective treo, tiêu chí
chấm thiếu hoặc không quan sát được), phủ luyện-tập-chủ-động cho objective bậc cao,
tính khả thi triển khai so với brief, và tính cụ thể của formative evaluation. Một
kiểm tra---\emph{measurability} của objective---được báo cho audit nhưng không tự
sửa, vì một Bloom verb quan sát được đã làm objective đo được và việc viết lại các
phát biểu vốn đã tốt sẽ làm giảm chất lượng (\cref{sec:discussion}).

\paragraph{Repair có chủ sở hữu.} Mỗi loại vi phạm thuộc về specialist mà artifact
của nó có thể sửa. Sau khi agent $A_i$ sinh, nó verify và repair chỉ các loại nó sở
hữu, tối đa một budget $B$:
\[
\text{lặp } \big[\, \Omega_i \leftarrow \{\omega \in \mathcal{V}(G) : \mathrm{owner}(\omega)=i\};\ G \leftarrow A_i^{\mathrm{repair}}(G, \Omega_i)\,\big]
\]
cho đến khi $\Omega_i = \emptyset$ hoặc hết budget. Một repair kết hợp một revision
nhắm-đích bằng LLM lên artifact bị đánh dấu với một fallback tất định đảm bảo hội tụ
ngay cả khi output của model không dùng được. Vì repair thuộc về chính specialist đã
sinh, việc bỏ decomposition (\cref{sec:experimental-methodology}) tất yếu bỏ luôn
specialist repair: một biến thể single-agent tự sửa toàn cục cho chính nó, nên không
specialist nào bị lén đưa trở lại.

\subsection{Aggregation và Residual Audit}
\label{sec:aggregation}

Một aggregator cuối kết xuất graph thành ADDIE output contract và kiểm nó theo
schema. Nó cũng ghi một \emph{residual audit im lặng}: chạy lại verifier một lần và
lưu mọi vi phạm còn lại làm metric, không sửa. Điều này phơi bày, với mỗi cấu hình,
bao nhiêu alignment defect còn sót---một thước đo điều kiện (i) độc lập với điểm
rubric dùng ở \cref{sec:results}.

\subsection{Thuật toán}
\label{sec:algorithm}

\begin{figure}[t]
\small
\hrule\vspace{2pt}
\textbf{Thuật toán 1:} Sinh với AlignmentGraph-ISD\vspace{2pt}
\hrule\vspace{4pt}
\begin{tabbing}
xx\=xx\=xx\=\kill
\textbf{Input:} brief $b$; specialist $A_1,\dots,A_k$ theo thứ tự ADDIE; budget $B$\\[1pt]
$G \gets \emptyset$\\
\textbf{for} $i = 1$ \textbf{to} $k$ \textbf{do}\\
\>$c \gets \rho(b) \cup \gamma(G) \cup \mu(M)$\\
\>$G \gets A_i(G, c)$ \quad // sinh; tạo node/edge sở hữu\\
\>\textbf{for} $r = 1$ \textbf{to} $B$ \textbf{do} \quad // self-check trong lượt\\
\>\>$\Omega_i \gets \{\omega \in \mathcal{V}(G) : \mathrm{owner}(\omega)=i\}$\\
\>\>\textbf{if} $\Omega_i = \emptyset$ \textbf{then break}\\
\>\>$G \gets A_i^{\mathrm{repair}}(G, \Omega_i)$\\
\textbf{return} $\mathrm{aggregate}(G)$, \ $\mathrm{residual}\gets\mathcal{V}(G)$
\end{tabbing}
\vspace{-4pt}\hrule
\caption{Sinh với self-check per-agent.}
\label{alg:harness}
\end{figure}

\Cref{alg:harness} tóm tắt việc sinh. Workflow là một pipeline cố định; không dùng
dynamic planner trong setting single-pass nghiên cứu ở đây (một biến thể dùng planner
được giữ lại cho phần mở rộng multi-session ở \cref{sec:threats-limitations}).

\subsection{Khác biệt so với các Baseline Gần nhất}
\label{sec:diff}

So với các baseline mạnh nhất, khác biệt mang tính biểu diễn chứ không chỉ quy trình.
ReAct-ADDIE và các agent bám lý thuyết khác truyền chỉ dẫn alignment qua prompt và
không bao giờ vật chất hoá nó; objective chúng sinh và assessment chúng sinh không
bao giờ được liên kết trong một cấu trúc mà hệ thống có thể kiểm tra.
AlignmentGraph-ISD thay vào đó biến mỗi quan hệ alignment thành một edge có kiểu được
tạo lúc sinh và được kiểm bởi một verifier có chủ sở hữu, nên khác biệt về input,
điều khiển và output là khác biệt về \emph{trạng thái nào tồn tại}: một graph mà
harness có thể cưỡng chế \cref{eq:aligned}, so với văn bản mà nó chỉ có thể prompt.
